%!TEX root = ../../../main.tex

\chapter{Vorverarbeitung und Suche}
\label{chap:vorverarbeitung}

\paragraph{Einlesen.} Der Parser liest PDF mit PyMuPDF und DOCX mit python-docx. Zuerst sucht er Tabellen: Jede Zeile wird ein eigener Block \enquote{Spalte: Wert | Spalte: Wert}, so steht der Spaltenkopf bei jeder Zahl. Überschriften erkennt er an Schriftgröße (ab 1,25-fach) oder Fettdruck und führt einen Kapitelpfad mit. Absätze unter 20 Zeichen hängt er an den vorigen Block, über 1500 Zeichen teilt er an Satzgrenzen. Bilder beschreibt gemma3:4b mit allen ablesbaren Werten. Aus 16 Dokumenten entstehen 436 Blöcke.

\begin{lstlisting}[language={}, basicstyle=\ttfamily\scriptsize, caption={Eine Tabellenzeile im JSON-Zwischenformat (gekürzt)}, label=lst:chunk, frame=single, breaklines=true]
{"chunk_type": "table_row", "content": "Gebührenart: Schmutzwassergebühr (Kanal) | Menge / Bemessung: 125,0 m³ | Satz: 2,80 €/m³ | ... | Brutto (€): 416,50",
 "heading_stack": ["ZWECKVERBAND WASSERVERSORGUNG", "Abgabenbescheid: Wasser & Abwasser 2024", "2. Gebührenberechnung ..."],
 "page_number": 1, "metadata": {"filename": "Rechnung_Wasser_Zweckverband_2024.pdf", "folder": "Rechnungen", "reconstructed": true}}
\end{lstlisting}

\paragraph{Speicherebenen und Kerndaten.} Ebene~1 ist eine lesbare JSON-Datei je Dokument (Listing~\ref{lst:chunk}). Ebene~2 speichert teure Modellausgaben unter einem SHA-256-Schlüssel, damit eine Parser-Änderung keine neuen Modellaufrufe kostet. Ebene~3 ist ChromaDB. Zusätzlich füllt qwen2.5:7b je Dokument 19 Kerndatenfelder. Code löscht jeden Wert, der nicht im Dokumenttext steht. 96\,\% der Felder stimmen mit meinen Sollwerten überein.

\paragraph{Suche.} Beim Einbetten (nomic-embed-text) setze ich vor jeden Block seinen Pfad \texttt{[Ordner | Datei | Kapitel]}. So wird \enquote{2,80\,€/m³} als Teil der \emph{Wasserabrechnung} gespeichert und nicht als irgendeine Zahl. ChromaDB liefert 20 Kandidaten nach Kosinus-Ähnlichkeit (\texttt{base}). Passt ein Wort der Frage zu Ordner, Datei oder Kapitel, wird mit 1,5 multipliziert (\texttt{boost}), und die besten fünf gehen an den Prototyp. BM25, einen Reranker und das Umformulieren der Frage habe ich gemessen (Anhang~\ref{app:messungen}). Keiner war besser, also sind sie aus.
